%=====================================================================
%  Persistent Memory Architecture in Agents Using DBMS
%  Phase 2 -- Comparison, Findings and Recommendation
%  Term Paper -- Thematic Assessment 2
%  Harshit Khemani (241302081), Kush Ahuja, Madhav Bassi,
%  Kushagra Agrawal
%  B.Tech CSE (AI/ML), Section C
%  Department of Computer Science & Engineering, SGT University
%
%  Typography, palette and rhythm carry over from Phase 1 so the two
%  reports read as one series. The content does not: no section of the
%  Phase 1 report is reproduced here.
%
%  Prose follows ASD-STE100: active voice, simple tense, short sentences.
%=====================================================================
\documentclass[11pt,a4paper]{article}

%---------------------------- fonts ---------------------------------
\usepackage[T1]{fontenc}
\usepackage[utf8]{inputenc}
\usepackage{XCharter}                    % body serif
\usepackage[scaled=0.95]{FiraSans}       % sans, for headings and labels
\usepackage[varqu,scaled=0.98]{zi4}      % Inconsolata mono
\usepackage[charter,vvarbb]{newtxmath}
\renewcommand{\familydefault}{\rmdefault}
\usepackage[letterspace=80]{microtype}

%---------------------------- layout --------------------------------
\usepackage[a4paper,top=2.8cm,bottom=2.8cm,left=3.75cm,right=3.75cm,
            headheight=14pt,footskip=1.1cm]{geometry}

\usepackage{graphicx}
\usepackage{booktabs}
\usepackage{tabularx}
\usepackage{array}
\usepackage[table,dvipsnames]{xcolor}
\usepackage{listings}
\usepackage{enumitem}
\usepackage{fancyhdr}
\usepackage{titlesec}
\usepackage{caption}
\usepackage{float}
\usepackage{tikz}
\usetikzlibrary{positioning,arrows.meta,fit,calc,shapes.geometric,
                decorations.pathreplacing}
\usepackage[hidelinks,breaklinks=true]{hyperref}
% Break a long URL at punctuation only. Breaking at any character, which
% is what xurl allows, splits "https:" from "//" and cuts words in half.
% hyperref rewrites these at \begin{document}, so set them after it. An
% empty \UrlBigBreaks stops the break after "https:", and a small
% \Urlmuskip keeps the break points from opening visible gaps.
\AtBeginDocument{%
  \def\UrlBreaks{\do\/\do\-\do\.\do\_\do\?\do\&\do\=}%
  \def\UrlBigBreaks{}%
  \Urlmuskip=0mu plus 1mu\relax}

\hypersetup{
  pdftitle={Persistent Memory Architecture in Agents Using DBMS ---
            Phase 2: Comparison, Findings and Recommendation},
  pdfauthor={Harshit Khemani (241302081), Kush Ahuja, Madhav Bassi,
             Kushagra Agrawal},
  pdfsubject={Database Management Systems --- Term Paper, Phase 2},
  pdfkeywords={agent memory, DBMS, constraints, request amplification,
               lost update, bitemporal, normalisation, PostgreSQL}
}

%---------------------------- colour --------------------------------
% The Phase 1 palette, unchanged. Blue carries structure. Clay carries
% anything the agent does not own. Every text colour clears 4.5:1.
\definecolor{ink}     {HTML}{1A1D21}
\definecolor{blue800} {HTML}{1E3A5F}
\definecolor{blue600} {HTML}{35618F}
\definecolor{blue200} {HTML}{C2D5E8}
\definecolor{blue100} {HTML}{E8EFF7}
\definecolor{clay700} {HTML}{8A5A2B}
\definecolor{clay600} {HTML}{A8703A}
\definecolor{clay200} {HTML}{E6D3BC}
\definecolor{clay100} {HTML}{F7EEE3}
\definecolor{grey600} {HTML}{5C6570}
\definecolor{grey300} {HTML}{CBD0D7}
\definecolor{grey100} {HTML}{F1F3F6}

\color{ink}

%------------------------ vertical rhythm ---------------------------
\linespread{1.16}\selectfont
\setlength{\parindent}{0pt}
\setlength{\parskip}{0.62\baselineskip}
\setlist{topsep=0.3\baselineskip,itemsep=0.18\baselineskip,parsep=0pt,
         leftmargin=1.25em}

\clubpenalty=10000
\widowpenalty=10000
\displaywidowpenalty=10000
\raggedbottom
\setlength{\emergencystretch}{1.5em}
\renewcommand{\arraystretch}{1.25}

% Each of the three criteria has a page budget, so floats have to pack
% tightly. LaTeX's defaults reserve a fifth of every page with a float
% for running text and refuse a top float larger than 70% of the page;
% both push a table to the next page and leave the current one short.
\renewcommand{\topfraction}{0.92}
\renewcommand{\bottomfraction}{0.75}
\renewcommand{\textfraction}{0.06}
\renewcommand{\floatpagefraction}{0.80}
\setcounter{topnumber}{3}
\setcounter{bottomnumber}{2}
\setcounter{totalnumber}{4}

%---------------------------- headers -------------------------------
\pagestyle{fancy}
\fancyhf{}
\fancyhead[L]{\footnotesize\sffamily\color{grey600}%
  Persistent Memory Architecture \textperiodcentered\ Phase 2}
\fancyhead[R]{\footnotesize\sffamily\color{grey600}%
  Khemani \textperiodcentered\ Ahuja \textperiodcentered\ Bassi
  \textperiodcentered\ Agrawal}
\fancyfoot[C]{\footnotesize\sffamily\color{grey600}\thepage}
\renewcommand{\headrulewidth}{0.4pt}
\renewcommand{\headrule}{\color{grey300}%
  \hrule height 0.4pt width \headwidth}
\fancypagestyle{plain}{\fancyhf{}%
  \fancyfoot[C]{\footnotesize\sffamily\color{grey600}\thepage}%
  \renewcommand{\headrulewidth}{0pt}}

%---------------------------- headings ------------------------------
\titleformat{\section}
  {\sffamily\bfseries\large\color{blue800}}{\thesection}{0.65em}{}
\titlespacing*{\section}{0pt}{1.5\baselineskip}{0.45\baselineskip}
\titleformat{\subsection}
  {\sffamily\bfseries\normalsize\color{ink}}{\thesubsection}{0.6em}{}
\titlespacing*{\subsection}{0pt}{1.05\baselineskip}{0.25\baselineskip}

%---------------------------- captions ------------------------------
\captionsetup{font={small},labelfont={sf,bf,color=blue800},
              labelsep=period,skip=7pt,width=0.95\textwidth}

%---------------------------- listings ------------------------------
\lstdefinestyle{base}{
  backgroundcolor=\color{grey100},
  basicstyle=\ttfamily\scriptsize\color{ink},
  keywordstyle=\color{blue800}\bfseries,
  stringstyle=\color{clay700},
  commentstyle=\color{grey600},
  numbers=none,
  showstringspaces=false,
  breaklines=true,
  frame=leftline,
  framerule=1.6pt,
  rulecolor=\color{blue200},
  framexleftmargin=6pt,
  xleftmargin=8pt,
  tabsize=2,
  captionpos=b,
  columns=fullflexible,
  keepspaces=true,
  aboveskip=0.7\baselineskip,
  belowskip=0.4\baselineskip
}
\lstset{style=base}
\lstdefinestyle{sql}{style=base,language=SQL,
  morekeywords={CONFLICT,GENERATED,ALWAYS,STORED,EXCLUDE,USING,GIST,
    RESTRICT,MATERIALIZED,CONCURRENTLY,tstzrange,int8range,RETURNING,
    REFERENCES,UNIQUE,DO,NOTHING,DATERANGE}}

%---------------------------- helpers -------------------------------
\newcommand{\code}[1]{\texttt{\small #1}}
% \code carries an absolute \small, which is correct in running text
% and too large inside a \scriptsize table. \ddl inherits the size of
% whatever cell it sits in, and its argument may carry its own \\.
\newcommand{\ddl}[1]{\texttt{\begin{tabular}{@{}l@{}}#1\end{tabular}}}
\newcommand{\tikzcode}[1]{\texttt{#1}}
\newcolumntype{L}[1]{>{\raggedright\arraybackslash}p{#1}}
\newcolumntype{R}[1]{>{\raggedleft\arraybackslash}p{#1}}
% A ragged X. Justified X columns this narrow open word-sized gaps.
\newcolumntype{Y}{>{\raggedright\arraybackslash}X}
% Fira's default figures are oldstyle, so a sans-set zero sits at
% x-height and "Mem0" reads as "Memo". This forces lining figures for
% one span without changing the numerals anywhere else.
\newcommand{\lnum}[1]{{\firalining #1}}
\newcommand{\khe}{\href{https://www.khe.money}{Harshit Khemani}}
\newcommand{\eyebrow}[1]{{\sffamily\footnotesize\bfseries\color{clay700}%
  \MakeUppercase{\textls[140]{#1}}}}
% A finding label. The recommendation cites these numbers by name.
\newcommand{\finding}[1]{{\sffamily\bfseries\color{blue800}#1}\quad}

%---------------------------- tikz ----------------------------------
\tikzset{every picture/.append style={
  execute at begin picture={\let\code\tikzcode\linespread{1.0}\selectfont}}}
\tikzset{
  box/.style   ={draw=blue200, line width=0.7pt, fill=blue100,
                 rounded corners=2pt, align=center, inner sep=5pt,
                 font=\sffamily\footnotesize},
  boxc/.style  ={draw=clay200, line width=0.7pt, fill=clay100,
                 rounded corners=2pt, align=center, inner sep=5pt,
                 font=\sffamily\footnotesize},
  boxn/.style  ={draw=grey300, line width=0.7pt, fill=white,
                 rounded corners=2pt, align=center, inner sep=5pt,
                 font=\sffamily\footnotesize},
  ar/.style    ={-{Stealth[length=2mm]}, draw=blue600, line width=0.7pt},
  arc/.style   ={-{Stealth[length=2mm]}, draw=clay600, line width=0.7pt},
  arn/.style   ={-{Stealth[length=2mm]}, draw=grey300, line width=0.7pt,
                 dashed},
  tag/.style   ={font=\sffamily\scriptsize, text=grey600, inner sep=2pt},
  card/.style  ={font=\sffamily\scriptsize, text=grey600, fill=white,
                 inner sep=1.5pt},
  note/.style  ={font=\sffamily\scriptsize, text=clay700, align=left},
  bar/.style   ={draw=none, fill=blue600},
  barc/.style  ={draw=none, fill=clay600},
  barl/.style  ={font=\sffamily\scriptsize, text=ink, anchor=east},
  barv/.style  ={font=\sffamily\scriptsize, text=grey600, anchor=west},
  gridline/.style={draw=grey300, line width=0.4pt}
}

\begin{document}

%=====================================================================
%  TITLE PAGE
%=====================================================================
\begin{titlepage}
\centering
\sffamily
\linespread{1.0}\selectfont
\setlength{\parskip}{0pt}
\vspace*{1.2cm}

{\footnotesize\bfseries\color{grey600}%
 \textls[140]{SHREE GURU GOBIND SINGH TRICENTENARY UNIVERSITY}}\\[3pt]
{\footnotesize\color{grey600}Gurugram, Haryana}\\[1.1cm]

{\color{blue200}\rule{\textwidth}{1pt}}\\[9pt]
{\large\bfseries\color{ink}Department of Computer Science \& Engineering}\\[7pt]
{\color{blue200}\rule{\textwidth}{0.5pt}}\\[1.6cm]

\eyebrow{Term Paper \textperiodcentered\ Thematic Assessment, Phase 2}\\[9pt]
{\normalsize\color{grey600}Database Management Systems Laboratory}\\[1.3cm]

% A tabular with a \strut in every row: each row then has identical
% height and depth, so the two 7pt gaps render exactly equal.
{\fontsize{26}{30}\selectfont\bfseries\color{blue800}%
\begin{tabular}{@{}c@{}}
\strut Persistent Memory\\[7pt]
\strut Architecture in\\[7pt]
\strut Agents Using DBMS
\end{tabular}}\\[14pt]

{\large\color{clay700}Comparison, Findings and Recommendation}\\[1.8cm]

\begin{minipage}{0.86\textwidth}
\small
\renewcommand{\arraystretch}{1.42}
\begin{tabular}{@{}r@{\hspace{1.5em}}l@{}}
{\color{grey600}Submitted by}   & \bfseries\khe\ \ \ {\mdseries 241302081} \\
                               & \bfseries Kush Ahuja \\
                               & \bfseries Madhav Bassi \\
                               & \bfseries Kushagra Agrawal \\[4pt]
{\color{grey600}Programme}     & B.Tech, Computer Science \& Engineering (AI/ML) \\
{\color{grey600}Section}       & C \\
{\color{grey600}Department}    & Computer Science \& Engineering \\
{\color{grey600}Institution}   & SGT University \\
{\color{grey600}Date of submission} & 10 September 2026 \\
\end{tabular}
\end{minipage}

\vfill
{\color{blue200}\rule{\textwidth}{0.5pt}}\\[6pt]
{\footnotesize\color{grey600}Academic Session 2026 \quad\textperiodcentered\quad
 Semester V \quad\textperiodcentered\quad Phase 2 of 2}
\end{titlepage}

%=====================================================================
%  FRONT MATTER
%=====================================================================
\pagenumbering{roman}
\setcounter{page}{2}

\section*{\sffamily\color{blue800}Abstract}
\addcontentsline{toc}{section}{Abstract}

Phase 1 of this term paper argued that the memory of an AI agent is a
database, and it described how four deployed systems keep facts across a
session boundary. Phase 2 tests that argument against measurement.

The subject of this phase is \textbf{Lore}, the open-source memory system
that Notion released for MCP-compatible assistants \cite{loreRepo,
notionBlog}. Lore is a good subject because it is honest about its own
trade. It stores agent memory as rows in five Notion databases, so the
memory stays readable by people and needs no server. In exchange it gives
up every guarantee a database declares: no unique index, no foreign key,
no check constraint, no atomic write across pages, no join, no aggregate
and no conditional write.

This report compares that design against two relational engines that run
the same schema, reports what the comparison found, and gives one
recommendation. The evidence comes from our companion study, which
reconstructed Lore's 80 properties from its source, generated a vault of
11\,885 memories and 4\,277 facts with definitional ground truth, and ran
a workload of 361 questions in ten classes through three stores
\cite{loreStudy}.

The result is not that the document store gives wrong answers. It gives
right ones. It answers 100.0\% of the workload exactly, and so does
SQLite; PostgreSQL answers 99.2\%, and its three divergences are a
stemming rule inside its text index rather than a fault. What separates
the three is \emph{cost} and \emph{refusal}. The same 361 questions cost the
Notion Data API 522\,686 requests against 361 SQL statements. Eight
agents writing one key through a read-then-write protocol lose 87.1\% of
their updates, and the same workload through a unique index loses none.
PostgreSQL declines even to add a temporal exclusion constraint to the
vault, because 70 subject--predicate pairs already contradict themselves.

The recommendation follows from those numbers, and it is not
``use PostgreSQL''. It is to keep the legible store as the system of
record and to add a rebuildable relational projection beside it, which
carries the indexes, the constraints and the joins the legible store
cannot declare.

\vspace{0.4\baselineskip}
\eyebrow{Keywords}\\[2pt]
agent memory, request amplification, lost update, referential integrity,
bitemporal data, exclusion constraint, normalisation, entity resolution,
PostgreSQL, SQLite.

\clearpage

{\setlength{\fboxsep}{9pt}%
\colorbox{clay100}{\parbox{\dimexpr\textwidth-18pt\relax}{\small
\eyebrow{What this report does not repeat}\\[3pt]
This is a standalone Phase 2 report. No section, table, figure or
argument from the Phase 1 report appears here. Phase 1 covered the four
storage layers, the PostgreSQL schema, the record and recall paths, and
four case studies. This phase covers only the three criteria set for it:
\textbf{comparison}, \textbf{findings} and \textbf{recommendation}. It
takes a new subject system, a new evidence base, and a bibliography that
shares no research paper with Phase 1. Where a Phase 1 conclusion is
needed, it is named in one sentence, never reproduced.}}}

\vspace{0.6\baselineskip}

\section*{\sffamily\color{blue800}New Research Added in This Phase}
\addcontentsline{toc}{section}{New Research Added in This Phase}

The brief for this phase asks for three more research papers. The three
below are new to this report, none of them appears in Phase 1, and each
does real work in the sections that follow rather than sitting in the
list. Table~\ref{tab:new} names what each one contributes and where it is
used.

\begin{table}[!htbp]
\centering
\caption[The three new research papers]{The three research papers added in Phase 2, and the work each
one does.}
\label{tab:new}
{\footnotesize
\renewcommand{\arraystretch}{1.3}
\begin{tabularx}{\textwidth}{L{3.3cm} Y L{1.5cm}}
\toprule
\sffamily\bfseries Paper & \sffamily\bfseries What it contributes here &
\sffamily\bfseries Used in \\
\midrule
\textbf{Zep}, a temporal knowledge graph for agent memory
\cite{zep,graphiti} &
Zep attaches a validity interval to every subject--predicate--object
edge and invalidates an old edge instead of deleting it. It is the
sharpest prior art for Lore's Facts relation, so it fixes the axis on
which the comparison of data models is fair. &
\S\ref{sec:cmp}, \S\ref{sec:f5} \\
\textbf{Mem0}, scalable long-term memory with explicit operations
\cite{mem0} &
Mem0 gives memory a small data manipulation language: the extractor
decides \code{ADD}, \code{UPDATE}, \code{DELETE} or \code{NOOP} for every
candidate. That names the axis Phase 1 left open, which is who owns the
update and where it is applied. &
\S\ref{sec:cmp}, \S\ref{sec:f7} \\
\textbf{LongMemEval}, a benchmark for long-term interactive memory
\cite{longmemeval} &
It isolates five abilities of a long-term memory, two of which ---
temporal reasoning and knowledge updates --- are exactly the two a
validity interval exists to serve, and it measures about a 30\% accuracy
drop for commercial assistants across sustained interactions. That ties
the validity interval to abilities a benchmark treats as core. &
\S\ref{sec:cmp}, \S\ref{sec:f5} \\
\bottomrule
\end{tabularx}}
\end{table}

Six further sources are also new to this report. Five support the
argument at specific points: Graphiti, the engine behind Zep
\cite{graphiti}, the SQL:2011 temporal standard \cite{sql2011}, Gray's
account of the lost update \cite{gray1976}, Swoosh on entity resolution
\cite{swoosh}, and a recent survey that maps the operations of an agent
memory onto the operations of a database \cite{du2025}. The sixth,
A-MEM \cite{amem}, states the opposing position on whether a memory
should have a schema at all.

\clearpage
{\sffamily\setlength{\parskip}{0pt}%
 \tableofcontents
 \vspace{1.4\baselineskip}
 \listoffigures
 \vspace{1.2\baselineskip}
 \listoftables}
\clearpage

\pagenumbering{arabic}
\setcounter{page}{1}

%=====================================================================
\section{Comparison}
\label{sec:cmp}
%=====================================================================

Phase 1 asked how an agent memory carries a fact across a session
boundary. That separates designs but does not price them. This phase asks
a narrower question with a harder answer. \emph{One schema, three stores:
what does each cost, and what will each refuse?}

\subsection{The subject, and where the numbers come from}

Lore is Notion's open-source memory system for MCP-compatible assistants
\cite{loreRepo,notionBlog}. A vault is one Notion page holding five
databases --- Projects, Topics, Memories, Entities and Facts --- all
reached through the Notion Data API. That API is the whole design: at
most 100 rows per request, 3 requests per second, and no join, no
aggregate, no unique index, no check constraint and no conditional write
\cite{notionApi}.

Every number here comes from our companion study \cite{loreStudy}, which
reconstructed Lore's 80 properties from \code{src/notion/schema.ts} at
commit \code{95c3558} and re-implemented the schema on SQLite and on
PostgreSQL. Its corpus is generated world-first, so every correct answer
is definitional: 11\,885 memories, 4\,277 facts of which 1\,196 are
closed, 1\,096 entities carrying 1\,153 aliases. A workload of 361
questions in ten classes then runs through three arms: an emulator
restricted to the documented Data API, and the two engines.

\subsection{Three agent memories, compared as data models}

A comparison needs a fair axis, and the three new papers supply it. Zep
attaches a validity interval to every edge of a temporal knowledge graph,
so an old claim is invalidated and not deleted \cite{zep,graphiti}. Mem0
gives the writer explicit operations, so a candidate becomes an
\code{ADD}, an \code{UPDATE}, a \code{DELETE} or a \code{NOOP}
\cite{mem0}. Lore sits between them, and Table~\ref{tab:designs} sets the
three side by side.

\begin{table}[!htbp]
\centering
\caption[Three memory designs]{Three agent memory designs as data models. The last row is the one
the rest of this report measures.}
\label{tab:designs}
{\scriptsize
\renewcommand{\arraystretch}{1.25}
\begin{tabularx}{\textwidth}{L{2.25cm} Y Y Y}
\toprule
 & \sffamily\bfseries Lore \cite{loreRepo} &
   \sffamily\bfseries Zep and Graphiti \cite{zep,graphiti} &
   \sffamily\bfseries \lnum{Mem0} \cite{mem0} \\
\midrule
Temporal model &
Bitemporal. Valid From, Valid Until, Observed At, Invalidated At &
Bitemporal. An edge records when it was true and when the graph learned
it &
Timestamps on the item. No validity interval \\

How an update happens &
The client reads, decides, then writes. Two requests, and a window
between them &
The engine closes the old edge and adds the new one &
An explicit operation chosen at write time: \code{ADD}, \code{UPDATE},
\code{DELETE} or \code{NOOP} \\

\midrule
\sffamily\bfseries The store can refuse a bad write &
\sffamily\bfseries No. Uniqueness is a convention in the client &
\sffamily\bfseries Yes, for what the engine declares &
\sffamily\bfseries Only as far as the writer decides \\
\bottomrule
\end{tabularx}}
\end{table}

One row is missing on purpose: only Lore is readable and correctable by a
person without a client, the trade \S\ref{sec:rec} defends. A-MEM argues
the opposite case to all three, that a fixed schema limits cross-task
adaptability \cite{amem}; this phase does not settle that, but none of
the failures in \S\ref{sec:findings} follows from having a schema. Each
follows from a store that cannot enforce one. LongMemEval justifies the
temporal row: two of the five abilities it isolates, temporal reasoning
and knowledge updates, are exactly the two a validity interval exists to
serve \cite{longmemeval}.

\subsection{One schema, three stores}

Table~\ref{tab:arms} is the centre of this comparison: the same schema,
the same 361 questions, three stores.

\begin{table}[!htbp]
\centering
\caption[The three arms]{The three arms on one workload of 361 questions. Requests and
client-side rows are means for one question. The rate-limit floor applies
Notion's documented 3 requests per second to the whole workload; it is a
lower bound on wall-clock, not a measurement of it, and no rate limit
applies to a local engine. The PostgreSQL arm is PGlite, the same source
tree and planner compiled to WebAssembly, so no throughput claim is made
for it.}
\label{tab:arms}
{\footnotesize
\renewcommand{\arraystretch}{1.28}
\begin{tabularx}{\textwidth}{X R{2.3cm} R{1.7cm} R{2.0cm}}
\toprule
\sffamily\bfseries Measure &
\sffamily\bfseries Lore on Notion &
\sffamily\bfseries SQLite &
\sffamily\bfseries PostgreSQL \\
\midrule
Exact-answer rate                   & 100.0\%   & 100.0\% & 99.2\% \\
\midrule
Requests, one question              & 1\,447.9  & 1.0     & 1.0 \\
Requests, whole workload            & 522\,686  & 361     & 361 \\
Rate-limit floor, whole workload    & 48.4\,h   & ---     & --- \\
\midrule
Share of the workload it can state  & 33.2\%    & 100\%   & 100\% \\
Invariants with no declarative form & 4 of 7    & 1 of 7  & 0 of 7 \\
\bottomrule
\end{tabularx}}
\end{table}

The last two rows need a note. The first counts what the store can
\emph{express}: 44.6\% of the questions need a join, 11.4\% an aggregate
and 11.1\% a transitive closure, so 66.8\% of the workload needs
something the Data API cannot state as a filter. The second counts what
it can \emph{refuse}: four of the seven invariants have no declarative
form on Notion and three of those are guarded by nothing at all, while
SQLite declares six, failing only on the temporal one for want of an
exclusion constraint, and PostgreSQL declares all seven.
Figure~\ref{fig:amp} in \S\ref{sec:f2} draws the uneven distribution
behind the 1\,447.9 mean.

So the comparison does not separate the three on correctness. It
separates them on cost, as a multiplier rather than a constant, and on
refusal, which is categorical: a database's central service is not
answering a question but declining to enter a state that makes a future
answer wrong.

%% The brief sets a page budget for each of the three criteria, so
%% each one starts on a fresh page and its span is countable.
\clearpage
%=====================================================================
\section{Findings}
\label{sec:findings}
%=====================================================================

Most of the findings below carry the name of an older database
phenomenon, because that is what they turn out to be.

\finding{F1} One result frames the rest. \textbf{Correctness is not where
these designs differ}, so nothing below is about a wrong answer. All
three stores return the exactly-correct set, and PostgreSQL's 99.2\% is
the stemming rule of \S\ref{sec:f8}, not a fault. Everything below is
about what the right answer costs and about what the store will not stop.

\subsection{A missing join costs a multiplier, not a constant}
\label{sec:f2}

\finding{F2} Provenance is the cheapest join an agent memory can be asked
for: a claim, the memory it came from, and that memory's author. In SQL
it is one statement. The Data API cannot filter one database by a
property of the page a relation points at, so the client pages over Facts
in 12 requests, reads 1\,126 memory pages one at a time and filters
itself --- 1\,137.75 round trips and 1\,126.2 client-side rows against
one and none, a floor of 6.3 minutes for one question.

Searching what the memories \emph{say} is worse. Property filters cannot
see page blocks and the search endpoint matches titles \cite{notionApi},
so a complete search is one list request per hundred rows plus one per
row: 11\,892 requests, 66.1 minutes here, 11.0 hours ten times larger,
110.1 hours a hundred times larger.

That figure carries the study's own biggest caveat: it measures what the
\emph{documented} API permits, and Lore delegates ranking to a Notion
endpoint the study could not measure, so Lore is cheaper here
\cite{loreStudy}. The structural point survives --- Lore owns no index,
so its retrieval depends on an endpoint it does not control.

\begin{figure}[!htbp]
\centering
\resizebox{\textwidth}{!}{%
\begin{tikzpicture}[x=1cm,y=1cm]
% Decade gridlines at 2.5 cm per decade: 1, 10, 100, 1000, 10000.
% 2.9 cm to the decade. Bar length is log10(value) times that, so the
% 1.0x class has no bar and 11 892x runs just past the last gridline.
\foreach \gx/\glab in {0/1,2.9/10,5.8/100,8.7/1000,11.6/10\,000}{
  \draw[gridline] (\gx,-0.14) -- (\gx,3.14);
  \node[font=\sffamily\scriptsize,text=grey600,anchor=north]
    at (\gx,-0.18) {\glab$\times$};
}
\foreach \nm/\ln/\vl/\yy/\st in {%
  body-search/11.818/11\,892$\times$/2.92/barc,
  provenance/8.863/1\,138$\times$/2.62/barc,
  conflict-scan/4.197/28.0$\times$/2.32/bar,
  cross-db/3.840/21.1$\times$/2.02/bar,
  supersession/1.624/3.6$\times$/1.72/bar,
  ask-entity/1.384/3.0$\times$/1.42/bar,
  as-of/1.384/3.0$\times$/1.12/bar,
  current/1.384/3.0$\times$/0.82/bar,
  aggregate/0.873/2.0$\times$/0.52/bar,
  wake-up/0/1.0$\times$/0.22/bar}{
  \node[barl] at (-0.22,\yy) {\nm};
  \fill[\st] (0,\yy-0.085) rectangle (\ln,\yy+0.085);
  \node[barv] at (\ln+0.14,\yy) {\vl};
}
\end{tikzpicture}}
\caption[Request amplification by question class]{Request amplification by question class: the mean number of Data
API requests that one question costs, against the single SQL statement
which answers it. The axis is logarithmic and evenly spaced by decade, so
\code{wake-up} at $1\times$ has no bar to draw. The two clay bars are the
two classes that pass a thousand requests for one question, and both are
searches the store cannot push down.}
\label{fig:amp}
\end{figure}

\finding{F2b} The prompt is bounded and the work behind it is not. The
wake-up hook caps its payload at 20 memories and 30 facts, and the cap
holds: over an eight-fold growth of the vault it still costs 24 requests
and injects between 775 and 852 tokens. What grows is the class of
operation touching every row --- the conflict scan, the debt sweep, the
aggregate, the complete search --- which a rate limit turns into hours.
Lore's own caps cluster there.

\subsection{A read-then-write upsert loses most of its writes}
\label{sec:f3}

\finding{F3} Lore upserts a topic-keyed memory by reading the vault for a
matching row and then creating or patching. The API documents no
conditional write --- no ETag, no \code{If-Match}, no version token
\cite{notionApi} --- so the window between read and write cannot be
closed from the client, and Lore's own source says so. The study ran the
interleave explicitly rather than leaving it to a scheduler: every writer
reads before any writer writes, the worst schedule the protocol admits
and the one the API gives no way to exclude. With 8 writers over 200
rounds the read-then-write path lost 1\,393 of 1\,600 updates, 87.1\%,
and left 7 duplicate rows for a key that should hold one. The same
workload through \code{INSERT \ldots ON CONFLICT} lost none.

The loss grows with contention. A shallower sweep at 2, 4, 8 and 16
writers lost 49.0\%, 73.5\%, 85.8\% and 91.9\%, against 0.0\% at every
width through the unique index. This is Gray's lost update
\cite{gray1976} by a new route: the anomaly needs no weak isolation level
when the protocol has no conditional write to close on.

\subsection{A constraint the store cannot declare is a state it cannot refuse}
\label{sec:f4}

\finding{F4} The study took seven invariants from the schema and wrote
the one line of standard DDL that would enforce each
(Table~\ref{tab:inv}). Four have no declarative form on the substrate at
all, and three of those are guarded by nothing whatsoever.

\begin{table}[!htbp]
\centering
\caption[The seven invariants]{Seven invariants the schema implies, the declaration that would
enforce each, and what walks through it today.}
\label{tab:inv}
{\scriptsize
\renewcommand{\arraystretch}{1.1}
\begin{tabularx}{\textwidth}{L{0.85cm} X L{3.05cm} L{2.5cm}}
\toprule
& \sffamily\bfseries Invariant & \sffamily\bfseries Declarative form &
  \sffamily\bfseries In the vault \\
\midrule
INV-1 & At most one memory for one topic key &
  \ddl{UNIQUE (topic\_key,\\ project\_id)} & Client code \\
INV-2 & Every fact points at a memory that exists &
  \ddl{FOREIGN KEY ...\\ ON DELETE RESTRICT} & \bfseries Nothing \\
INV-3 & One object per functional predicate at any instant &
  \ddl{EXCLUDE USING gist\\ (... WITH \&\&)} & \bfseries Nothing \\
INV-4 & Confidence is one of three values &
  \ddl{CHECK (confidence\\ IN ...)} & Editable options \\
INV-5 & An alias names exactly one entity &
  \ddl{UNIQUE (alias)} & \bfseries Nothing \\
INV-6 & Valid from is not after valid until &
  \ddl{CHECK (... <= ...)} & Not verifiable \\
INV-7 & The dedup key is the hash of its own columns &
  \ddl{GENERATED ALWAYS\\ AS (...)} & Write-time code \\
\bottomrule
\end{tabularx}}
\end{table}

Three were attacked. INV-1 is F3 above; in the other two the vault
refused nothing at all. Over 2\,724 facts on functional predicates,
PostgreSQL under a temporal exclusion constraint accepted 2\,532 and
refused 192, and the vault refused 0. When 112 memories were archived
after facts already cited them, \code{ON DELETE RESTRICT} refused all 112
and left no orphan; the vault accepted all 112 and left 112 claims whose
provenance it can no longer return. The exclusion result has a sharper
first half: the constraint could not be added to the vault \emph{as
generated} at all, because 70 subject--predicate pairs already assert two
objects over overlapping time.

Lore is not unaware of contradictions --- it ships a scanner, capped at
500 candidates for each project. The comparison is between preventing a
state and detecting it afterwards at quadratic cost under a cap.
Detection is useful, and it is not the same thing.

\subsection{Two time axes disagree often enough to be a correctness decision}
\label{sec:f5}

\finding{F5} A bitemporal schema is only worth its columns if the two
axes ever disagree. Here they disagree constantly. A retraction lags the
world by a median of 23 days and a 95th percentile of 43, and of the
1\,196 facts whose validity has closed, 1\,172 --- 98.0\% --- were
retracted only after the claim had stopped being true. A fact disagrees
on day $D$ when ``was this true on $D$'' and ``did the vault believe it
on $D$'' give different answers; across all 4\,277 facts, 245 disagree on
day 120 (5.7\%), 123 on day 240 (2.9\%), 112 on day 360 (2.6\%) and 55 on
day 480 (1.3\%).

Lore and Zep both hold the two axes \cite{loreRepo,zep}, and SQL:2011
standardised the separation more than a decade before either
\cite{sql2011}. LongMemEval measures what getting it wrong costs:
assistants lose about 30\% of their accuracy across sustained
interactions \cite{longmemeval}. A query that does not name its axis is a
correctness error waiting for a probe day.

\subsection{The normalisation violations that ship are substrate costs}
\label{sec:f6}

\finding{F6} The usual exercise walks a badly shaped relation up to
Boyce--Codd form. The interesting exercise here is the other one: which
violations \emph{ship}, and why. Five rungs of a six-rung ladder are
still in the released product, and none is an oversight. Aliases are a
comma-joined list in one text cell, so no predicate can ask which entity
owns an alias; a multi-select would force a schema migration for every
new alias. Two derived columns stand in for one expression index, because
the store has none. And the centre of the knowledge graph carries a
Boyce--Codd violation that is structural rather than careless: every
Notion database needs exactly one title property, a title cannot be a
relation, so a fact holds its subject twice.

That violation is reachable by the most ordinary maintenance the system
has. The merge routine repoints every fact relation from loser to winner
and never rewrites the title, and the vault was generated to hold what
that leaves behind: 105 facts, 2.5\% of all, whose title says one name
and whose relation points at another. Entity resolution has an older
theory --- Swoosh defines the transitive closure a resolution must reach
\cite{swoosh} --- and Lore's merge is careful but single-hop, so 39
subjects have their facts split across two rows.

\textbf{One measurement went the other way.} The study predicted that
resolving an alias through the joined cell would be expensive, because a
substring test over-matches, and built 57 aliases as deliberate prefixes
of others. Across 1\,057 alias terms the filter returned 1\,159 rows of
which 1\,153 were exact: six discarded, 0.5\% of the traffic. The
magnitude did not support the argument, so it is reported as measured.

\subsection{The update anomaly lands in the prose, not in the table}
\label{sec:f7}

\finding{F7} This is the finding a schema alone does not fix. A fact is
restated across several memories in prose, and when the fact changes the
Fact row is updated and the prose is not. Across the 1\,196 closed facts,
3\,867 memories still assert the closed value: a mean of 3.23 stale
restatements each, and for 33 facts the stale copies alone would fill a
ten-row recall. An agent that retrieves memories rather than facts
retrieves the old answer, in a confident sentence, nothing marking it
retired.

This is the classical update anomaly. Mem0 answers it by making
\code{UPDATE} and \code{DELETE} explicit operations a writer must choose
\cite{mem0}. Lore supersedes the fact row, which is correct and
incomplete, because the copies live where supersession does not reach.

\subsection{Three indexes over one text give two different answers}
\label{sec:f8}

\finding{F8} The workload turned up a result it was not designed to find.
Seven test phrases returned identical counts on all three indexes ---
\emph{call graph} matched 677 memories on each, \emph{error rate} 200,
\emph{deploy manifest} 152. The eighth did not. \emph{Timed out} appears
nowhere as a literal string, because the memories say \emph{times out},
so a literal search and SQLite's FTS5 both return 0 and PostgreSQL
returns 669: the English configuration stems \emph{timed} to \emph{time}
and drops \emph{out} as a stop word, so the phrase collapses to the
lexeme \code{\textquotesingle time\textquotesingle} and matches every
memory that mentions timing at all.

None of the three is wrong --- a stemmer is a considered position on what
``discusses this'' means. What matters is that the position is a choice
and no answer states which choice was made, so two teams on two engines
reason from two different sets of facts.

%% The brief sets a page budget for each of the three criteria, so
%% each one starts on a fresh page and its span is countable.
\clearpage
%=====================================================================
\section{Recommendation}
\label{sec:rec}
%=====================================================================

The recommendation is not ``use PostgreSQL''. Moving off the document
substrate forfeits three things that are not conveniences: a surface
non-engineers can audit, an authorisation model already provisioned, and
no operational estate at all.

\textbf{Keep the legible store as the system of record, and put a
relational projection in front of it.} The projection holds the indexes,
the constraints and the joins the substrate cannot declare, and holds no
authority: it derives from the vault, so if it is wrong, drop it and
derive it again. Because a constraint cannot refuse a write it never
sees, the write paths differ. \textbf{The agent writes through the
projection}, where a bad write is refused before it reaches the vault,
which is what makes rules 1 and 3 true rather than aspirational.
\textbf{People still write straight to the vault}, so a human edit that
breaks an invariant fails the next rebuild instead --- detection rather
than prevention, on a schedule nobody forgets. \textbf{Every read goes to
the projection}, where a join, an aggregate and a transitive closure are
each one statement.

Six rules follow. Each names the finding forcing it; F1 forces none,
being the framing result rather than a defect.

\begin{enumerate}
  \item \textbf{Declare the key where the data is.} A unique index and
        one \code{ON CONFLICT} statement close the window a
        read-then-write upsert cannot; the same index on an alias stops
        one real thing being introduced twice. \emph{(F3, F4, F6)}
  \item \textbf{Put the join on the engine's side of the wire.} Move the
        classes needing a join, an aggregate or a transitive closure ---
        two thirds of this workload. \emph{(F2)}
  \item \textbf{Give the fact table a temporal primary key.} An exclusion
        constraint over a range type refuses the contradiction, leaving
        the scan nothing to find \cite{sql2011}. \emph{(F4, F5)}
  \item \textbf{Make every query name its axis.} Was the claim true on
        day $D$, or did the store believe it on $D$? \emph{(F5)}
  \item \textbf{Retire the restatements when you retire the fact.} Give
        the extractor an explicit \code{UPDATE} and \code{DELETE}, as
        Mem0 does \cite{mem0}. \emph{(F7)}
  \item \textbf{Record what answered, and log what a cap dropped.} A
        silent truncation reads as complete coverage. \emph{(F2b, F8)}
\end{enumerate}

\textbf{When, and with what caveat.} Not on day one, but where
maintenance passes stop finishing inside a session or more than one agent
writes one key. The measurements bound the shape of the cost rather than
any team's wall-clock, coming from an emulator over a generated corpus
\cite{loreStudy}. What they settle is firmer: a memory that arrives at
triples with validity intervals, an entity registry, provenance edges and
a confidence score has arrived at a database schema, and the only
question left is which guarantees it does without, and what it does
instead.

%=====================================================================
\section*{\sffamily\color{blue800}References}
\addcontentsline{toc}{section}{References}

Entries \textbf{[1]--[3]} are the three research papers added in this
phase. Entries \textbf{[4]--[7]} are the subject system and the evidence
base. Entries \textbf{[8]--[13]} support the argument and are also new to
this report. No entry is carried over from Phase 1.

% thebibliography opens with a \section*{\refname} of its own, and this
% report has already set that heading. Swallow the duplicate: the two
% arguments taken here are the star and the \refname group. Nothing after
% the bibliography uses \section.
\renewcommand{\section}[2]{}
\begin{thebibliography}{99}
% Ragged right here. Justifying a reference list forces a long URL to
% break at a bad point rather than move whole to the next line.
{\small\raggedright\setlength{\itemsep}{0.35\baselineskip}

\bibitem{zep} P. Rasmussen, P. Paliychuk, T. Beauvais, J. Ryan and
D. Chalef. ``Zep: A Temporal Knowledge Graph Architecture for Agent
Memory.'' arXiv:2501.13956, 2025.
\url{https://arxiv.org/abs/2501.13956}

\bibitem{mem0} P. Chhikara, D. Khant, S. Aryan, T. Singh and D. Yadav.
``Mem0: Building Production-Ready AI Agents with Scalable Long-Term
Memory.'' \emph{ECAI 2025}, Frontiers in Artificial Intelligence and
Applications, vol. 413, pp. 2993--3000, 2025; arXiv:2504.19413.
\url{https://arxiv.org/abs/2504.19413}

\bibitem{longmemeval} D. Wu, H. Wang, W. Yu, Y. Zhang, K.-W. Chang and
D. Yu. ``LongMemEval: Benchmarking Chat Assistants on Long-Term
Interactive Memory.'' \emph{ICLR 2025}; arXiv:2410.10813.
\url{https://arxiv.org/abs/2410.10813}


\bibitem{loreRepo} Notion Labs, Inc. (makenotion). ``Lore --- persistent,
shared AI memory backed by Notion for MCP-compatible assistants.''
GitHub, MIT licence, read at commit \code{95c3558}.
\url{https://github.com/makenotion/lore}

\bibitem{notionBlog} H. Salman. ``Building Shared Memory for AI Agents in
Notion.'' Notion Blog, 18 August 2026.
\href{https://www.notion.com/blog/building-shared-memory-for-ai-agents-in-notion}%
{\nolinkurl{notion.com/blog/building-shared-memory-for-ai-agents-in-notion}}

\bibitem{notionApi} Notion Labs, Inc. ``Request Limits'', ``Pagination''
and ``Search by Title.'' Notion API Reference.
\url{https://developers.notion.com/reference/request-limits}

\bibitem{loreStudy} H. Khemani, K. Ahuja, M. Bassi and K. Agrawal.
``Persistent Memory Architecture for Agents: A Review of Notion's Lore.''
Review paper, measured 22 August 2026.
\url{https://dbms-memory.khe.money}


\bibitem{graphiti} Zep AI. ``Graphiti --- build temporal context graphs
for AI agents.'' GitHub, Apache-2.0 licence.
\url{https://github.com/getzep/graphiti}

\bibitem{sql2011} K. Kulkarni and J.-E. Michels. ``Temporal Features in
SQL:2011.'' \emph{ACM SIGMOD Record}, 41(3):34--43, 2012.

\bibitem{gray1976} J. N. Gray, R. A. Lorie, G. R. Putzolu and
I. L. Traiger. ``Granularity of Locks and Degrees of Consistency in a
Shared Data Base.'' In \emph{Modelling in Data Base Management Systems},
North Holland, 1976, pp. 365--394.

\bibitem{swoosh} O. Benjelloun, H. Garcia-Molina, D. Menestrina, Q. Su,
S. E. Whang and J. Widom. ``Swoosh: A Generic Approach to Entity
Resolution.'' \emph{The VLDB Journal}, 18(1):255--276, 2009.

\bibitem{amem} W. Xu, Z. Liang, K. Mei, H. Gao, J. Tan and Y. Zhang.
``A-MEM: Agentic Memory for LLM Agents.'' arXiv:2502.12110, 2025.
\url{https://arxiv.org/abs/2502.12110}

\bibitem{du2025} Y. Du, W. Huang, D. Zheng, Z. Wang, S. Montella,
M. Lapata, K.-F. Wong and J. Z. Pan. ``Rethinking Memory in AI: Taxonomy,
Operations, Topics, and Future Directions.'' arXiv:2505.00675, 2025.
\url{https://arxiv.org/abs/2505.00675}
}
\end{thebibliography}

\vspace{1.2\baselineskip}
{\color{blue200}\rule{\textwidth}{0.5pt}}\\[5pt]
{\footnotesize\sffamily\color{grey600}
\khe\ (241302081) \quad\textperiodcentered\quad Kush Ahuja
\quad\textperiodcentered\quad Madhav Bassi \quad\textperiodcentered\quad
Kushagra Agrawal\\[2pt]
B.Tech CSE (AI/ML), Section C \quad\textperiodcentered\quad
Department of Computer Science \& Engineering, SGT University\\[2pt]
Phase 2 \quad\textperiodcentered\quad submitted 10 September 2026}

\end{document}
